% Ch18 WS1 -- galvanic cell anatomy and operation. Block 1.
\documentclass{apchem-worksheet}

\apcunitword{Chapter}
\apcunit{18}
\apcunittitle{Electrochemistry}
\apcdoctitle{Worksheet 1 \textbullet\ Galvanic Cells and How They Work}
\apcsubtitle{Zumdahl \S18.1 \textbullet\ An Ox, Red Cat --- true in every cell there is}

\begin{document}
\wsheader[Never label an electrode positive or negative --- the AP exam
excludes it. Use \textbf{anode} and \textbf{cathode}. $E^\circ$ (V):
\ce{Ag+/Ag} $+0.80$ \textbullet\ \ce{Cu^2+/Cu} $+0.34$ \textbullet\
\ce{Pb^2+/Pb} $-0.13$ \textbullet\ \ce{Fe^2+/Fe} $-0.44$ \textbullet\
\ce{Zn^2+/Zn} $-0.76$ \textbullet\ \ce{Mg^2+/Mg} $-2.37$.]

\problem[]{Name the component of a galvanic cell responsible for each
function:}
\begin{parts}
  \item Site of oxidation: \answerblank[4.4cm]{the anode}
  \item Site of reduction: \answerblank[4.4cm]{the cathode}
  \item Carries electrons between electrodes:
        \answerblank[4.4cm]{the external wire}
  \item Keeps each half-cell electrically neutral:
        \answerblank[4.4cm]{the salt bridge}
  \item Measures the driving force:
        \answerblank[4.4cm]{the voltmeter}
\end{parts}

\problem[]{Explain in detail why a galvanic cell stops almost immediately if
the salt bridge is removed.
\answerlines[4]{As the cell runs, cations enter solution at the anode and
are removed at the cathode, so the anode compartment builds up positive
charge and the cathode compartment builds up negative charge. That charge
separation opposes further electron flow through the wire and halts it. The
salt bridge lets spectator ions migrate --- anions toward the anode, cations
toward the cathode --- cancelling the imbalance without letting the two
solutions mix.}}

\problem[]{Why does separating the two half-reactions produce useful energy
when simply dropping zinc into copper(II) sulfate does not?
\answerlines[3]{In direct contact the electrons pass straight from zinc
atoms to copper ions and the energy is released as heat. Separating the
half-cells forces the electrons to travel through an external wire to get
from one to the other, and that flow is an electric current, which can do
work.}}

\problem[]{Consider $\ce{Zn(s) + Cu^2+(aq) -> Zn^2+(aq) + Cu(s)}$ run as a
galvanic cell.}
\begin{parts}
  \item Write the two half-reactions and label each.
        \answerlines[2]{Anode (oxidation): \ce{Zn -> Zn^2+ + 2e^-}.
        Cathode (reduction): \ce{Cu^2+ + 2e^- -> Cu}.}
  \item Which electrode loses mass, and why?
        \answerlines[2]{The zinc. Zinc atoms are oxidized to \ce{Zn^2+} and
        leave the electrode to enter the solution.}
  \item Which electrode gains mass, and why?
        \answerlines[2]{The copper. \ce{Cu^2+} ions are reduced to copper
        atoms, which deposit onto the electrode.}
  \item In which direction do electrons flow through the wire?
        \answerblank[7.6cm]{from zinc (anode) to copper (cathode)}
  \item In which direction do anions move in the salt bridge?
        \answerblank[6cm]{toward the anode}
  \item What happens to $[\ce{Cu^2+}]$ as the cell runs?
        \answerblank[4cm]{it decreases}
\end{parts}

\problem[]{Complete the comparison:}
\begin{parts}
  \item Galvanic cells involve a reaction that is
        \answerblank[5.4cm]{thermodynamically favored}
  \item Electrolytic cells involve a reaction that is
        \answerblank[6.2cm]{thermodynamically unfavored}
  \item Galvanic cells have $E^\circ_{\text{cell}}$
        \answerblank[2.4cm]{positive}, electrolytic cells
        \answerblank[2.4cm]{negative}
  \item In \emph{both} types, oxidation occurs at the
        \answerblank[2.6cm]{anode} and reduction at the
        \answerblank[2.6cm]{cathode}
  \item Which type requires an external power source?
        \answerblank[3cm]{electrolytic}
\end{parts}

\problem[]{A student builds a cell from magnesium and copper. After an hour
the magnesium strip is visibly thinner and the copper strip is coated.}
\begin{parts}
  \item Identify the anode and cathode.
        \answerblank[6cm]{anode Mg, cathode Cu}
  \item Write the overall reaction.
        \answerblank[6.4cm]{\ce{Mg + Cu^2+ -> Mg^2+ + Cu}}
  \item Calculate $E^\circ_{\text{cell}}$. \workspace[1.8cm]
        \keyonly{\begin{apcanswer}$E^\circ = 0.34 - (-2.37) =
        \mathbf{+2.71}~\text{V}$\end{apcanswer}}
  \item Is the cell galvanic or electrolytic? Justify from your answer.
        \answerlines[2]{Galvanic --- the cell potential is positive, so the
        reaction is thermodynamically favored and the cell produces
        electrical energy rather than consuming it.}
\end{parts}

\problem[]{A student writes: ``The anode is the negative electrode, so
electrons flow away from it.'' Explain why the second half of that sentence
is fine but the first half should not appear on an AP answer.
\answerlines[3]{Electrons do flow away from the anode --- that is correct
and always true, because oxidation there releases them. But the sign of an
electrode is not something the AP exam assesses: the convention reverses
between galvanic and electrolytic cells, which is exactly why the CED
excludes it. Anode and cathode are defined by oxidation and reduction and
never change meaning.}}

\problem[]{Predict, with a calculation, whether each reaction occurs
spontaneously as written:}
\begin{parts}
  \item \ce{Pb(s) + 2Ag+(aq) -> Pb^2+(aq) + 2Ag(s)} \workspace[1.8cm]
        \keyonly{\begin{apcanswer}$E^\circ = 0.80 - (-0.13) = +0.93$~V ---
        \textbf{yes}, favored\end{apcanswer}}
  \item \ce{Cu(s) + Zn^2+(aq) -> Cu^2+(aq) + Zn(s)} \workspace[1.8cm]
        \keyonly{\begin{apcanswer}$E^\circ = -0.76 - 0.34 = -1.10$~V ---
        \textbf{no}, unfavored; this is the reverse Daniell
        cell\end{apcanswer}}
  \item \ce{Fe(s) + Cu^2+(aq) -> Fe^2+(aq) + Cu(s)}
        \answerblank[4.4cm]{$+0.78$ V; yes}
\end{parts}

\begin{spiralstrip}{Chapter 17 \textbullet\ spiral}
  \item Favored means $\Delta G^\circ$ is: \answerblank[2cm]{negative}
  \item Crossover temperature $=$ \answerblank[2.6cm]{$\Delta H/\Delta S$}
  \item $\Delta S^\circ$ sign when gas moles rise:
        \answerblank[1.8cm]{positive}
  \item Favored but very slow means: \answerblank[3.4cm]{kinetic control}
  \item $S^\circ$ for an element is: \answerblank[2cm]{not zero}
\end{spiralstrip}

\end{document}
